\chapter{Low Latency Engineering}
\label{chap:lowlatency}

Earlier chapters give the building blocks: pinning and scheduling 
(Chapter~\ref{chap:linuxprocesses}), caches and false sharing 
(Chapter~\ref{chap:cacheperformance}), NUMA (Chapter~\ref{chap:numa}), 
non-blocking I/O and \texttt{epoll} (Chapters~\ref{chap:linuxnetworking} 
and~\ref{chap:eventdrivenio}), and kernel bypass when the stack itself is the limit 
(Chapter~\ref{chap:kernelbypass}). 
This chapter is the \highlight{ops checklist}: isolate cores, poll instead of sleep, 
keep the hot path cache-friendly and allocation-free, and optimize for latency even when 
that hurts bulk throughput.

\section{CPU Isolation}
\label{sec:cpu-isolation}
Assumes per-thread pinning from Section~\ref{subsec:sched-setaffinity}. 
This section covers kernel boot parameters and IRQ routing that isolate cores 
from general scheduler and interrupt noise.
\newline
Pinning alone is not enough if the kernel still steals your core for timers, RCU callbacks, 
hard IRQs, or unrelated threads. Isolation makes a core \highlight{mostly yours}: 
each knob below removes a different class of jitter. Use them as a \highlight{set}, 
not as independent toggles.

\subsection{\texttt{isolcpus}: scheduling exclusion}
\label{subsec:isolcpus}
By default the Linux scheduler treats every online CPU as a shared worker: 
the CFS load balancer can migrate runnable threads onto any core that looks idle. 
\texttt{isolcpus=}~\cite{kernel-isolcpus} removes listed CPUs from that regular 
balancing pool (exact flags and semantics depend on kernel version). 
Those cores no longer absorb random housekeeping or application threads unless 
you explicitly place work there with \texttt{sched\char`_setaffinity}.
\newline
\textbf{Noise removed:} the balancer treating your core as free capacity for 
unrelated tasks.
\newline
Typical pattern: housekeeping (IRQs, OS daemons, logging) on CPUs 0--$N$; 
strategy / market-data threads pinned onto isolated CPUs $N{+}1$ onward. 
\texttt{isolcpus} alone is incomplete --- timer ticks, RCU callbacks, and hard 
IRQs can still hit the core (sections below).

\subsection{\texttt{nohz\_full}: adaptive-ticks / tickless}
\label{subsec:nohz-full}
The kernel normally arms a \highlight{periodic timer interrupt} on each CPU 
(the scheduling ``tick'', historically every few milliseconds) for accounting, 
timers, and preemption decisions. Even a single pinned busy-loop thread gets 
interrupted by that tick --- pure jitter on a latency path.
\newline
\texttt{nohz\char`_full=}~\cite{kernel-nohz} enables \highlight{adaptive-ticks} 
(also called tickless / full dynticks) on the listed CPUs: when only one runnable 
task is present, the kernel can stop the periodic tick and wake the CPU mainly 
for real work. Needs a kernel built with full dynticks support; validate on your distro.
\newline
\textbf{Noise removed:} the periodic scheduling-clock interrupt that would otherwise 
preempt your pinned thread even when nothing else is runnable.
\newline
Pair with \texttt{isolcpus} and \emph{one} pinned busy-loop or poll thread per core; 
multiple runnable tasks on the same CPU defeat the point.

\subsection{\texttt{rcu\_nocbs}: RCU callback offload}
\label{subsec:rcu-nocbs}
\textbf{RCU} (Read-Copy-Update) is a Linux synchronization style for read-heavy 
shared data: readers proceed almost without locking; writers publish a new version 
and \emph{defer} freeing the old one until every concurrent reader is done. 
That deferred cleanup runs as \highlight{RCU callbacks} (often via softirq / 
dedicated kthreads).
\newline
By default, callback work can run on the same CPUs that queued it --- including 
an ``isolated'' core --- and inject softirq latency into your hot path. 
\texttt{rcu\char`_nocbs=}~\cite{kernel-rcu} (and related 
\texttt{rcu\char`_nocb\char`_poll}) marks listed CPUs as \emph{no-callback} CPUs 
and moves that processing onto housekeeping cores.
\newline
\textbf{Noise removed:} deferred RCU cleanup / softirq work on the cores you 
wanted quiet.
\newline
Together with the other knobs: \texttt{isolcpus} + \texttt{nohz\char`_full} + 
\texttt{rcu\char`_nocbs} + IRQ affinity + pinned threads.

\subsection{\texttt{irqaffinity} / \texttt{smp\_affinity}: hard-IRQ routing}
\label{subsec:irq-affinity}
A \highlight{hard IRQ} is a device interrupt (NIC queue, timer, disk controller, \ldots) 
delivered by hardware to some CPU. Unlike a voluntary context switch, it 
\emph{forces} that CPU into kernel interrupt context --- disastrous for 
tick-to-trade tails if the NIC's IRQs land on your strategy core.
\newline
\texttt{isolcpus} does not automatically move IRQs. You must steer them onto 
housekeeping CPUs~\cite{kernel-irq-affinity}:
\begin{itemize}
    \item \texttt{irqaffinity=} (boot): default CPU mask for IRQs that do not set 
          their own affinity --- keep that default on housekeeping cores.
    \item \texttt{/proc/irq/<n>/smp\char`_affinity} (or 
          \texttt{smp\char`_affinity\char`_list}): per-IRQ mask; pin each NIC 
          queue IRQ explicitly away from isolated strategy cores.
    \item \texttt{irqbalance}: userspace daemon that reshuffles IRQ affinity at 
          runtime for throughput fairness. On latency boxes, disable it or 
          constrain it --- otherwise it undoes your masks.
\end{itemize}
\par\noindent\textbf{Interview tip:} say the full set --- scheduler exclusion, 
tickless, RCU offload, hard-IRQ routing, then pinned threads --- not just 
``we set \texttt{isolcpus}''.

\section{Busy Polling}
\label{sec:busy-polling}
Builds on non-blocking sockets (Section~\ref{sec:blocking-vs-nonblocking-io}) and \texttt{epoll} 
(Chapter~\ref{chap:eventdrivenio}). Trades CPU usage for lower wake-up latency.
\newline
Blocking \texttt{epoll\char`_wait} / \texttt{recv} sleeps the thread; wake-up adds 
unpredictable delay. Alternatives on the critical path:
\begin{itemize}
    \item \texttt{epoll\char`_wait} with timeout \texttt{0} in a tight loop (spin until ready).
    \item Socket \highlight{busy poll} options (\texttt{SO\char`_BUSY\char`_POLL} and related 
          sysctls) so \texttt{recv}/\texttt{poll} spin briefly in the kernel networking stack 
          before sleeping~\cite{man-busy-poll}.
    \item User-space poll loops on bypass rings (Chapter~\ref{chap:kernelbypass}).
\end{itemize}
Cost: a core at $\sim$100\% for as long as you spin. Only do this on isolated CPUs dedicated 
to that job.
\par\noindent\textbf{Interview tip:} sleep = jitter; busy poll = latency for watts and opportunity cost.

\section{Cache Optimization}
\label{sec:lowlatency-cache-optimization}
See Chapter~\ref{chap:cacheperformance} for cache hierarchy, false sharing, and 
\texttt{perf}-based profiling. This section covers operational cache tuning only.
\begin{itemize}
    \item Keep hot structures within few cache lines; align and pad to avoid false sharing.
    \item Prefer sequential access and preallocated contiguous buffers (rings, arenas).
    \item Pin thread and memory on the same NUMA node as the NIC (Chapter~\ref{chap:numa}).
    \item Warm critical code and data at startup (touch pages, run a dry path) to avoid 
          cold-cache and major-fault surprises (Chapter~\ref{chap:memorymapping}).
    \item Measure with \texttt{perf} --- do not guess which miss hurts the tick-to-trade path.
\end{itemize}

\section{Avoiding Allocations}
\label{sec:avoiding-allocations}
See Section~\ref{sec:spsc-ring-buffer} for pre-allocated SPSC ring buffers. 
This section covers general allocation avoidance on hot paths.
\newline
\texttt{new}/\texttt{malloc}, \texttt{std::string} growth, and hidden temporaries cause:
\begin{itemize}
    \item Unbounded latency (allocator locks, syscalls, page faults),
    \item Cache/TLB disruption,
    \item Jitter under load that never shows in a quiet unit test.
\end{itemize}
Practices:
\begin{itemize}
    \item Preallocate pools and rings at startup; recycle objects 
          (Section~\ref{subsec:memory-pools}).
    \item Fixed-size messages; no unbounded buffering on the hot thread.
    \item Defer logging, formatting, and allocations to a colder thread via a lock-free queue.
    \item Prefer stack/static storage or arenas with LIFO free for short-lived data.
\end{itemize}

\section{Latency vs Throughput}
\label{sec:latency-vs-throughput}
\begin{center}
\begin{tabular}{|l|p{5.8cm}|p{5.8cm}|}
\hline
\textbf{Goal} & \textbf{Optimize for} & \textbf{Often sacrifices} \\
\hline
Latency & Median and tail response time; 
  predictable wake-ups; short paths & CPU utilization, batching efficiency, 
  multi-tenant density \\
\hline
Throughput & Messages/sec or bandwidth; 
  batching, coalescing, deep queues & Per-message delay; tail latency under load \\
\hline
\end{tabular}
\end{center}
HFT market-data and order paths are usually \highlight{latency-first}: busy poll, 
\texttt{TCP\char`_NODELAY}, shallow queues, one connection per hot thread, isolation. 
A bulk file transfer or analytics job is throughput-first: large buffers, interrupts 
coalesced, sleep when idle.
\par\noindent\textbf{Interview tip:} know which axis you are on. Techniques that maximize packets/sec 
(Nagle, interrupt coalescing, shared noisy cores) often \highlight{hurt} tick-to-trade 
latency --- and the reverse wastes capacity if you only care about bulk bandwidth.
